\begin{proposition}[The Liu--Rollenske family and its transverse deformation]\hypertarget{tag-01GL}{}\label{cdp:st:E14}
\StatementRecord{cdp:st:E14}{proved}{01GL}{\SeriesRef{cdp:eq:E14-elliptic-family}, \SeriesRef{cdp:st:E13}, \SeriesRef{cdp:eq:E14-canonical}, \SeriesRef{cdp:eq:E14-jet-map}, \SeriesRef{cdp:eq:E14-log-coho}, \SeriesRef{cdp:st:E6}, \SeriesRef{cdp:st:E1}; \cite[Proposition~2.1]{LiuRollenske2026}; \cite[Theorem~5.1]{LiuRollenske2026}}

Every member $Y(u)$ of the Liu--Rollenske family, with its scale
sufficiently small and nonzero as in its construction, satisfies
\[
 h^\bullet(Y(u),T_{Y(u)})=(1,2,2,0).
\]
Its Kuranishi germ is a smooth reduced surface. One parameter is
the given scale; a second parameter replaces its cuspidal fibre by
two rank-one nodal fibres. All sufficiently small deformations
retain a torus fibration and have algebraic dimension one.
\end{proposition}


\begin{proof}

\begin{enumerate}[label=\arabic*.,leftmargin=*]
\item In the Weierstrass equation of \cite[Proposition~2.1]{LiuRollenske2026}, make the explicit
coordinate comparison
\[
 t_{\rm new}=\frac{27}{4}(t_{\rm old}^{-1}-1),\quad
 x_{\rm new}=x_{\rm old}/t_{\rm old}^2,\quad
 y_{\rm new}=y_{\rm old}/t_{\rm old}^3.
\]
Put $k=27/4$. Near $t_{\rm old}=0$, the new infinity coordinates are
$x_{\rm old}/[k(1-t_{\rm old})]^2$ and
$y_{\rm old}/[k(1-t_{\rm old})]^3$, and its base coordinate is
$t_{\rm old}/[k(1-t_{\rm old})]$. Near $t_{\rm old}=\infty$, put
$v=t_{\rm old}^{-1}$; the new fibre coordinates are
$v^2x_{\rm old},v^3y_{\rm old}$ and the new base coordinate is $k(v-1)$.
These charts extend the displayed maps to an isomorphism of the
Weierstrass models and their minimal resolutions, preserving $O$ and $P$.
They identify $\mathcal O(P-O)$ and its lifted translation; using the
transported linearization preserves the scale $u$.
The new equation is $y^2=x^3+tx+t$, with section $P=(-1,i)$; the multiple
value is now infinity, the cusp is at zero and the remaining nodal
value is $-27/4$. Consider the explicit deformation
\begin{equation}\hypertarget{tag-01A8}{}\label{cdp:eq:E14-elliptic-family}
 y^2=x^3+tx+t+\epsilon,\qquad
 P_\epsilon=(-1,b),\quad b^2=\epsilon-1,
 \quad b(0)=i.
\end{equation}
Its finite discriminant is $4t^3+27(t+\epsilon)^2$, whose discriminant
as a cubic in $t$ is $314928\epsilon^3(1-\epsilon)$.
The nodal value near $-27/4$ remains simple and holomorphically
labelled. The other two nodes coalesce to the type-II cusp at
$\epsilon=0$. Infinity keeps type $III^*$; on its order-four
good-reduction cover the equation is
\[
 y'^2=x'^3+x'+r^2+\epsilon r^6,
 \qquad P'_{\epsilon}=(-r^2,br^3).
\]

\item On a disc containing the cusp
and its two nearby nodal values, the total elliptic surface in
\eqref{cdp:eq:E14-elliptic-family} is smooth over $\epsilon$:
$\partial(y^2-x^3-tx-t-\epsilon)/\partial t=-(x+1)$ is a unit
near the cusp. Translation by $-2P_\epsilon$ is an automorphism
there. Put $M_\epsilon=\mathcal O(P_\epsilon-O)$ and let $c(\epsilon)$
be the simple nodal value continuing $-27/4$. The identity
$\mathcal Hom(t_{2P}^*M,M)\simeq\mathcal O(F)$ in
\cite[Proposition~2.1]{LiuRollenske2026} persists in the smooth family
of elliptic surfaces: $H^1(\mathcal O)=H^2(\mathcal O)=0$ centrally
and nearby, so the exponential sequence identifies their Picard
groups with the fixed integral $H^2$ lattice. The two line bundles
have the same Chern class throughout the family. Their ratio is
trivial on every surface; Grauert's theorem makes this trivialization
holomorphic in $\epsilon$. The section of $\mathcal O(F_{c(\epsilon)})$
with divisor $F_{c(\epsilon)}$ therefore gives a holomorphic lifted
translation, agreeing with the original one at $\epsilon=0$.
For a uniformly small nonzero scale its contraction on the
complement of the zero section is uniform on the compact elliptic
fibres. The quotient is consequently free, properly discontinuous,
proper over the disc, and smooth over the parameter. This quotient
itself fills the cusp and the nearby rank-one nodal fibres. At the
separate node the ordinary rank-two Mumford charts vary
holomorphically. At infinity the fixed order-four torsion translation
and good-reduction comparison vary holomorphically in the displayed
equation. Gluing along their marked punctured collars therefore
constructs a proper holomorphic family over $(\epsilon,u)$ whose
central scale slice is $Y(u)$. Every member has the central
differentiable type by Ehresmann. The topology theorem
\cite[Theorem~5.1]{LiuRollenske2026} thus
identifies these differentiable manifolds with $S^6$.

\item Let $B$ be the reduced multiple fibre, so $f^*(\infty)=4B$.
The line $\mathcal O_Y(B)$, its canonical section and the isomorphism
$\mathcal O_Y(4B)\simeq f^*\mathcal O_{\mathbb P^1}(\infty)$ define
the root factorization $\bar f:Y\to\mathfrak B=\mathbb P(1,4)$.
For the coarse projection $\rho:\mathfrak B\to\mathbb P^1$ one has
$f=\rho\bar f$ and
$\rho^*\mathcal O_{\mathbb P^1}(1)=\mathcal O_{\mathfrak B}(4)$.
On the central threefold the
canonical form on the contracting quotient is the elliptic-surface
two-form wedged with $dv/v$. At the cusp it extends without a zero
or pole: from $t=(y^2-x^3)/(1+x)$ one gets
\[
 dt\wedge\frac{dx}{2y}=-\frac{dx\wedge dy}{1+x}.
\]
In a hexagonal chart $xyz=t-c(\epsilon)$, the form
$d\log x\wedge d\log y\wedge dt$ is $dx\wedge dy\wedge dz$.
At infinity the substitution
$t=r^{-4}$, $x=r^{-2}x'$, $y=r^{-3}y'$ gives the factor $-4r^{-4}$
in the surface form. The good three-form is invariant under the
free order-four affine action and descends to a nonvanishing frame.
These are the local canonical maps used in \ref{cdp:st:E13}.
Since the elliptic surface has
$K_S=\pi^*\mathcal O_{\mathbb P^1}(-1)$, these comparisons give
\begin{equation}\hypertarget{tag-01A7}{}\label{cdp:eq:E14-canonical}
 K_Y=f^*\mathcal O_{\mathbb P^1}(-1)
     =\bar f^*\mathcal O_{\mathfrak B}(-4).
\end{equation}
The calculation is simultaneous in $\epsilon$. Adjunction on every
full fibre gives trivial dualizing sheaf. \cite[Theorem~2.2, pp.~4--5]{Viaclovsky2026Combs} gives $h^0=1$;
duality gives $h^2=1$; constancy of its Euler characteristic in the
proper flat family gives $h^1=2$. Thus, also on the two-dimensional
base $(t,\epsilon)$, Grauert's theorem gives locally free $R^qf_*\mathcal O$
with base change, for $q=0,1,2$.

\item The multiplicative connection
\[
 \alpha=\frac{y+b}{x+1}\frac{dx}{2y},\qquad
 \nabla(hs)=(dh+\alpha h)s
\]
has precisely the residues of the rational section $s$ with divisor
$P_\epsilon-O$. For a vector $V$ dual to $dx/(2y)$, the horizontal lift is
$\widetilde V=V-\alpha(V)v\partial_v$. The multiplicative
connection commutes with the lifted translation, by the normalized
group-law comparison in \ref{cdp:st:E13}. On the good-reduction
chart its relative form is
\[
 \alpha=\frac{y'+br^3}{x'+r^2}\frac{dx'}{2y'},\qquad
 \frac{dx}{2y}=r\,\frac{dx'}{2y'}.
\]
Together with the scalar vector $v\partial_v$, this transition gives
the logarithmic Lie bundle $E=\mathcal O\oplus\mathcal O(-1)$. There the vector dual to $dx/(2y)$ is the derivation
$2y\partial_x+3x^2\partial_y$ on the cuspidal fibre, and its
horizontal lift is well defined because
$(y+b)/(x+1)$ is regular near the singular locus.

\item Put $G=R^1\bar f_*\mathcal O$. The inverse pairing associated with the invariant alternating
class gives $G\simeq E$ over the two-dimensional $(t,\epsilon)$
base away from its single cuspidal point: at a smooth fibre it is the pairing
$H^1(F,\mathbb C)/F^1\to(F^1)^*$; at a nodal fibre its quotient
frame consists of the graph-loop class and the elliptic normalization
class; and the good-reduction cover gives the frame at infinity.
These local pairing and normalization maps are the ones in
\ref{cdp:st:E13}, applied to the displayed connection and root transition. Both $G$ and $E$ are locally free across that point, by the
preceding cohomological-flatness and connection calculations.
The matrix entries of the isomorphism and its inverse extend
across that codimension-two set by Hartogs' theorem. Their products
are the identity on a dense open set, hence everywhere. Thus
\[
 G=\mathcal O\oplus\mathcal O(-1),\qquad
 R^2\bar f_*\mathcal O=\mathcal O(-1)
\]
on the central root curve as well; the second equality also follows
directly from relative duality and \eqref{cdp:eq:E14-canonical}.

\item Let $D=C+W$ be the sum of the cuspidal and hexagonal fibres.
For the local tangent calculation choose coordinates in which
$t=y^2-x^3$. They are obtained from the earlier formula by the
invertible change $\widetilde x=x/(1+x)^{1/3}$,
$\widetilde y=y/(1+x)^{1/2}$, with the branches equal to one at
zero, and then renaming the two new coordinates only for this local
logarithmic-module calculation. The logarithmic
tangent module is freely generated by
\[
 2y\partial_x+3x^2\partial_y,\quad
 2x\partial_x+3y\partial_y,\quad\partial_z,
\]
where $z$ is a coordinate along the elliptic singular curve. The
second generator maps to $6t\partial_t$;
the first and third generate the relative logarithmic kernel.
Consequently the global exact sequence is
\[
 0\longrightarrow\bar f^*E\longrightarrow T_Y(-\log D)
 \longrightarrow\bar f^*\mathcal O(-3)\longrightarrow0,
\]
because $T_{\mathfrak B}(-C-W)=\mathcal O(5-8)$.
The vertical cohomology is $(1,1,0,0)$. Every line in the
horizontal Leray terms is $\mathcal O(-3)$ or $\mathcal O(-4)$,
which has no cohomology on $\mathbb P(1,4)$. Hence
\begin{equation}\hypertarget{tag-01AA}{}\label{cdp:eq:E14-log-coho}
 h^\bullet(T_Y(-\log D))=(1,1,0,0).
\end{equation}

\item Return to the Weierstrass coordinates. The preceding coordinate
change preserves the ideal $(x^2,y)$ and induces the identity on
$\mathbb C[x]/(x^2)$, since $\widetilde x\equiv x\bmod x^2$.
The quotient at $W$ is $\eta_*K_{dP_6}$ with cohomology $(0,0,1)$.
At $C$, the map $v\mapsto v(t)|_C$ identifies the quotient with
the Jacobian ideal
\[
 J=(x^2,y)\subset\mathcal O_C.
\]
Indeed its kernel is precisely the logarithmic tangent module and
$N_{C/Y}=\mathcal O_C$. Put $Z=V(J)\subset C$ and consider
the restriction map
\begin{equation}\hypertarget{tag-01A9}{}\label{cdp:eq:E14-jet-map}
 H^1(C,\mathcal O_C)\longrightarrow H^1(Z,\mathcal O_Z).
\end{equation}

\item In the deformation \eqref{cdp:eq:E14-elliptic-family} the critical
base near the cusp is the smooth curve
\[
 \Gamma:\quad t=-3x^2,\qquad
 \epsilon=x^2(3+2x).
\]
Its map to the $\epsilon$-disc is finite flat of degree two; its
central fibre is $\operatorname{Spec}\mathbb C[x]/(x^2)$.
The lifted translation fixes its $x$-coordinate. This follows on
the punctured curve because translation fixes the unique node of
each nodal cubic, and then follows at zero by holomorphic
continuation. On the line coordinate it is
$(x,v)\mapsto(x,c(x)v)$ with $c$ a holomorphic nowhere-zero
contraction. Thus its quotient critical locus is exactly the
elliptic family
\[
 \mathcal Z=(\Gamma\times\mathbb C^*)/
       \langle(x,v)\mapsto(x,c(x)v)\rangle.
\]
In particular $H^0(Z,\mathcal O_Z)$ and
$H^1(Z,\mathcal O_Z)$ are both free rank-one modules over
$A=\mathbb C[x]/(x^2)$, and $H^2(Z,\mathcal O_Z)=0$.

\item Restriction in this proper family is a holomorphic bundle map
$G|_\Gamma\to R^1(\mathcal Z/\Gamma)_*\mathcal O_{\mathcal Z}$.
At a node with coordinate $x$, the toric Lie line is generated by
$\widetilde V+\frac{b}{1+x}v\partial_v$: the horizontal lift has
line-coordinate component $-\alpha(V)=-b/(1+x)$ at that node.
The quotient to the elliptic Lie line, and hence the normalization
restriction under the pairing above, therefore sends
$(v\partial_v,\widetilde V)$ to
$(1,-b/(1+x))$, up to a common nonzero factor. This is the local
map in \ref{cdp:st:E13}. Thus on $\Gamma\setminus\{0\}$ restriction,
after a nonzero scalar change of target frame, has the row
\[
 \left(1,-\frac{b}{1+x}\right).
\]
The normalization of $C$ is a $\mathbb P^1$-bundle over
the elliptic singular curve $E_0$, and its normalization quotient
is $\mathcal O_{E_0}$: the first jet at the cusp transforms with
derivative one under translation on the additive smooth locus.
The normalization exact sequence, together with
$h^1(\mathcal O_C)=2$, $h^2(\mathcal O_C)=1$ and
$h^\bullet(\mathcal O_{\widetilde C})=(1,1,0)$, implies
that $H^1(\mathcal O_C)\to H^1(\mathcal O_{\widetilde C})$
is onto. Restriction from the ruled normalization to its singular
section identifies the latter group with $H^1(\mathcal O_{E_0})$.
Thus restriction to $E_0$ is onto. The common scalar in the row
above cannot have a zero at zero, since both entries would then
vanish and this restriction would be zero. It is a unit.

\item The equality of these holomorphic maps on the punctured curve
extends over all of $\Gamma$, and in particular over its
nonreduced central fibre. Hence \eqref{cdp:eq:E14-jet-map} is,
in the specified bases, the complex-linear map
\[
 \mathbb C^2\longrightarrow A,\qquad
 (a,d)\longmapsto a+d(-i+ix).
\]
Its matrix in $1,x$ has determinant $i$, so it is an isomorphism.

\item The exact sequence $0\to J\to\mathcal O_C\to\mathcal O_Z\to0$
now gives
\[
 h^\bullet(C,J)=(0,1,1).
\]
Together with \eqref{cdp:eq:E14-log-coho} and the hexagonal quotient,
the tangent exact sequence gives
$h^\bullet(T_Y)=(1,2,2,0)$.

\item In the above local trivialization
of the total space the map to the base is
\[
 f_{\epsilon}=\frac{y^2-x^3-\epsilon}{1+x}.
\]
The derivative $\partial_\epsilon f$ has image $-1+x$ in
$\mathcal O_C/J=\mathbb C[x]/(x^2)$ along the singular curve.
If the ambient first-order deformation were trivial, every
first-order variation of its map would be a global section of
$f^*T_{\mathbb P^1}$, hence the pullback of a projective vector
field, because $f_*\mathcal O_Y=\mathcal O_{\mathbb P^1}$.
Its local image in this Jacobian quotient is constant. Changing
the local ambient trivialization adds $df(v)$ and hence adds
zero in that quotient. The nonzero coefficient of $x$ therefore
proves that the $\epsilon$-direction is nontrivial and cannot be
cancelled by the scale direction, which does not change this local
map. With $\epsilon=0$, a normalized singular elliptic curve is
$\mathbb C^*/\langle c_0u\rangle$; its modulus has nonzero
$u$-derivative for sufficiently small $u$. A trivial deformation
of the total space preserves its anticanonical pencil from
\eqref{cdp:eq:E14-canonical}, its three distinguished special values,
and that elliptic singular curve. Hence the scale direction is
nontrivial too. The two directions span $H^1(T_Y)$.
Present the Kuranishi local ring as
$\mathbb C\{z_1,z_2\}/I$. The classifying map of the constructed
bidisc, followed by the embedding in $H^1(T_Y)$, has invertible
differential and is a local analytic isomorphism by the inverse
function theorem. Every element of $I$ pulls back to zero, so $I=0$.
This is the ideal argument of \ref{cdp:st:E13}, and it proves that
the displayed family is locally miniversal over the full reduced germ.
The rank-two hypothesis in \ref{cdp:st:E6} is supplied for this
Liu--Rollenske construction by \ref{cdp:st:E7}. Thus all sufficiently
small deformations retain their torus fibration.
Their second Betti number remains zero, and \ref{cdp:st:E1}
identifies their meromorphic function field with the pullback of
$\mathcal M(\mathbb P^1)$, giving algebraic dimension one.

\end{enumerate}
\end{proof}
